% !TeX root = ../../../../../main.tex
% chapters/sections/chapter3/subsections/representative_household/lop.tex

\subsection{一物一価の法則}
\label{sec:chapter3_representative_household_lop}

\subsubsection{自国財の一物一価}
\label{sec:chapter3_representative_household_lop_home}

\paragraph{①自国の元となる方程式}
\label{sec:chapter3_representative_household_lop_home_original}
式
\eqref{form:chapter3_assumptions_lop_home_p_h_eq}
で仮定されたように、自国家計\(h\)の財の価格\(p_t^h\)について一物一価の法則は以下のようであった。
\begin{equation}
\label{form:chapter3_representative_household_lop_home_original_p_h_eq}
    p_t^h=e_t^{/*}p_t^{h*}
\end{equation}

\paragraph{②自国代表的家計の方程式}
\label{sec:chapter3_representative_household_lop_home_representative}
自国の生産者物価指数（PPI）\(p_t^H\)は式
\eqref{form:chapter3_households_stage1a_ppi_descriptions_p_H_eq}
によって示されたとおり、以下のように表現された。
\begin{equation}
\label{form:chapter3_representative_household_lop_home_representative_p_H_eq}
    p_t^H=\left[\sum_{h\in H}(p_t^h)^{1-\theta^H}\right]^{\frac{1}{1-\theta^H}}
\end{equation}
この右辺に先ほどの一物一価の法則
\eqref{form:chapter3_representative_household_lop_home_original_p_h_eq}
を代入すると以下のようになる。
\begin{equation}
\label{form:chapter3_representative_household_lop_home_representative_p_h_modified_eq}
    \begin{aligned}
        p_t^H&=\left[\sum_{h\in H}(e_t^{/*}p_t^{h*})^{1-\theta^H}\right]^{\frac{1}{1-\theta^H}} \\
        &=e_t^{/*}\left[\sum_{h\in H}(p_t^{h*})^{1-\theta^H}\right]^{\frac{1}{1-\theta^H}}
    \end{aligned}
\end{equation}
右辺の括弧は外国通貨建ての自国生産者物価指数\(p_t^{H*}\)の式
\eqref{form:chapter3_households_stage1a_ppi_descriptions_p_H_star_eq}
の右辺そのものであるから、代入すると\(p_t^{H*}\)に書き換えることができる。
こうして自国代表的家計の一物一価の法則が得られる。
\begin{equation}
\label{form:chapter3_representative_household_lop_home_representative_p_H_eq}
    p_t^H=e_t^{/*}p_t^{H*}
\end{equation}

\subsubsection{外国財の一物一価}
\label{sec:chapter3_representative_household_lop_foreign}

\paragraph{①外国の元となる方程式}
\label{sec:chapter3_representative_household_lop_foreign_original}
同様に式
\eqref{form:chapter3_assumptions_lop_foreign_p_f_eq}
で仮定されたように、外国家計\(f\)の財の価格\(p_t^f\)について一物一価の法則は以下のようであった。
\begin{equation}
\label{form:chapter3_representative_household_lop_foreign_original_p_f_eq}
    p_t^f=e_t^{/*}p_t^{f*}
\end{equation}

\paragraph{②外国代表的家計の方程式}
\label{sec:chapter3_representative_household_lop_foreign_representative}
外国の生産者物価指数（PPI）\(p_t^F\)は式
\eqref{form:chapter3_households_stage1a_ppi_descriptions_p_F_eq}
によって示されたとおり、以下のように表現された。
\begin{equation}
\label{form:chapter3_representative_household_lop_foreign_representative_p_F_def_eq}
    p_t^F=\left[\sum_{f\in F}(p_t^f)^{1-\theta^F}\right]^{\frac{1}{1-\theta^F}}
\end{equation}
この右辺に先ほどの一物一価の法則
\eqref{form:chapter3_representative_household_lop_foreign_original_p_f_eq}
を代入すると以下のようになる。
\begin{equation}
   \begin{aligned}
\label{form:chapter3_representative_household_lop_foreign_representative_p_f_modified_eq}
    p_t^F&=\left[\sum_{f\in F}(e_t^{/*}p_t^{f*})^{1-\theta^F}\right]^{\frac{1}{1-\theta^F}} \\
    &=e_t^{/*}\left[\sum_{f\in F}(p_t^{f*})^{1-\theta^F}\right]^{\frac{1}{1-\theta^F}}
   \end{aligned}
\end{equation}
右辺の括弧は外国生産者物価指数\(p_t^{F*}\)の式
\eqref{form:chapter3_households_stage1a_ppi_descriptions_p_F_star_eq}
の右辺そのものであるから、代入すると\(p_t^{F*}\)に書き換えることができる。
こうして外国代表的家計の一物一価の法則が得られる。
\begin{equation}
\label{form:chapter3_representative_household_lop_foreign_representative_p_F_lop_eq}
    p_t^F=e_t^{/*}p_t^{F*}
\end{equation}